\documentclass[11pt]{article}
\usepackage[UTF8]{ctex}
\usepackage{amsmath,amssymb,amsfonts,mathtools}
\usepackage{braket}
\usepackage[margin=1in]{geometry}
\usepackage{hyperref}
\usepackage{booktabs}
\hypersetup{colorlinks=true,linkcolor=blue,urlcolor=blue}

\title{\textbf{Problem 2 --- 选做 3}\\EWF Bath Threshold: MP2 自然轨道占据数与截断阈值}
\date{}
\begin{document}\maketitle

\section*{题目回顾}
MP2 bath natural orbitals (BNO) 的占据数大致按指数衰减; 截断阈值 $\eta$ 控制 bath
轨道数。给定 H$_2$O, $\eta\in[10^{-6},10^{-2}]$。要求:
\begin{enumerate}
  \item 推导 $n_{\text{BNO}}$ 与 $\eta$ 的渐近关系;
  \item 画出 fragment size 与能量误差随 $\eta$ 的变化, 找拐点, 解释
        ``diminishing returns'' (边际收益递减)。
\end{enumerate}

%=============================================================================
\section{背景: 为什么需要 MP2 BNO bath?}
%=============================================================================
\textbf{DMET (Schmidt) minimal bath} 用 HF 密度矩阵 $\gamma^{\text{HF}}$
的幂等性 ($(\gamma^{\text{HF}})^2{=}\gamma^{\text{HF}}$) 得到 bath 上界
$n_{\text{bath}}^{\text{DMET}}\le n_{\text{imp}}$, 它\emph{精确} 再现 HF 在
impurity 上的约化密度矩阵, 但只捕捉\emph{静态相关} (近简并、键断裂的纠缠)。

\textbf{当求解器升级到相关波函数 (CCSD/SQD/FCI)}, 相关态的密度矩阵
$\gamma^{\text{corr}}$ 不再幂等, 谱中出现的\emph{分数占据} ($0<n_k<1$) 反映
电子瞬时涨落 (动态相关)。这些涨落贡献了大量 ``超出 DMET 秩'' 的 Schmidt 分量,
DMET minimal bath 全部截掉 $\to$ fragment solver 看不到完整的动态相关空间 $\to$
引入 bath 截断误差 $\epsilon_{\text{frag}}$。

\textbf{EWF 的补救}: 在 DMET minimal bath 之外, \emph{额外} 加上由 MP2 密度修正
$\Delta\gamma^{\text{MP2}}$ 对角化得到的 bath natural orbitals (BNO), 按占据数
$\{n_i\}$ 截断: 保留 $n_i\ge\eta$ 的那些。MP2 是动态相关的最便宜近似, 故 BNO 是
``廉价捕动态相关'' 的最优 bath 选择。本题研究阈值 $\eta$ 与 bath 规模 / 精度
的标度关系。

%=============================================================================
\section{MP2 BNO 占据数 $n_i$ 的构造}
%=============================================================================
\subsection*{从 MP2 T2 振幅到密度修正}
MP2 波函数相对 HF 的一阶微扰 (in Rayleigh--Schr\"odinger 意义) 由
T2 振幅 $t_{ij}^{ab}$ 描述 ($i,j$ 占据, $a,b$ 虚):
\[
  |\Psi^{\text{MP2}}\rangle = |\Phi_0\rangle + \frac{1}{4}\sum_{ijab} t_{ij}^{ab}\,
  a^\dagger_a a^\dagger_b a_j a_i\,|\Phi_0\rangle .
\]
其密度矩阵修正 (相对 HF 占据投影)
\[
  \Delta\gamma^{\text{MP2}}
  = \gamma^{\text{MP2}}-\gamma^{\text{HF}}
\]
由 T2 振幅的二次型给出 (pyscf 标准实现):
\begin{align}
  \Delta\gamma_{ij}  &= 2\sum_{kab} t_{ik}^{ab}\,(t_{jk}^{ab})^*
                    - \sum_{kab} t_{ik}^{ab}\,(t_{jk}^{ba})^*  \quad \text{(占据块, 取负)},\\
  \Delta\gamma_{ab}  &= 2\sum_{ijc} t_{ij}^{ac}\,(t_{ij}^{bc})^*
                    - \sum_{ijc} t_{ij}^{ac}\,(t_{ij}^{cb})^*  \quad \text{(虚块, 取正)}.
\end{align}
MP2 是\emph{弱相关} 近似, 故 $\|\Delta\gamma^{\text{MP2}}\|\ll 1$; 其环境-环境块
$\Delta\gamma^{\text{MP2}}_{\text{env,env}}$ 是\emph{小半正定矩阵} (占据块和虚块
各自的二次型 $XX^\dagger$ 形式), 描述环境轨道如何被 ``瞬时双激发'' 激发到与
fragment+DMET-bath 子空间耦合的状态。

\subsection*{BNO = $\Delta\gamma^{\text{MP2}}_{\text{env,env}}$ 的本征轨道}
\textbf{bath natural orbital} 定义为 $\Delta\gamma^{\text{MP2}}_{\text{env,env}}$
的本征向量:
\[
  \Delta\gamma^{\text{MP2}}_{\text{env,env}}\,|b_i\rangle = n_i\,|b_i\rangle,
  \qquad n_1\ge n_2\ge\cdots\ge 0.
\]
对应本征值 $\{n_i\}$ (称为 \emph{BNO 占据数}) 反映: 若把环境轨道 $|b_i\rangle$
补进 bath, 它能给 fragment cluster 带来多少 MP2 相关能 / 多少密度涨落耦合。
占据数大 $\to$ 该 BNO 与 fragment 有强相关耦合 (重要);
占据数小 $\to$ 该 BNO 几乎不被激发 (可忽略)。
\textbf{截断阈值 $\eta$}: 仅保留 $n_i\ge\eta$ 的 BNO 作为 bath, 丢掉其余。

%=============================================================================
\section{$n_{\text{BNO}}$ vs $\eta$ 的渐近关系}
%=============================================================================
\subsection*{(1) 指数衰减假设}
BNO 占据数对应 ``双激发到环境轨道的概率'', 在弱相关体系 (H$_2$O 平衡几何) 中,
第 $i$ 个最相关 BNO 大致按\emph{指数}衰减:
\[
  \boxed{\;n_i \;\sim\; n_1\,e^{-\alpha\,(i-1)},\qquad \alpha>0\;}
  \tag{E}
\]
($\alpha$ 是衰减常数, 依赖体系/基组; H$_2$O/$6\text{-}31$G 实测
$\alpha\approx 2$--$5$ per BNO index, 见 \texttt{solution.py} 拟合。)
\textbf{物理起源}: 弱相关下 T2 振幅 $t_{ij}^{ab}$ 随激发轨道 $a,b$ 的轨道能差
$\sim 1/(\epsilon_a+\epsilon_b-\epsilon_i-\epsilon_j)$ 快速下降, 密度修正
$\Delta\gamma\sim t^2$ 是\emph{平方} 效应, 双重求和后大致按指数 (而非幂律)
衰减; 这是 RPA/MP2 density fitting 等近似下 known 的 ``double-excitation
weight'' 衰减规律。

\subsection*{(2) 对数截断关系 (核心结论)}
保留条件 $n_i\ge\eta$ 代入 (E):
\[
  n_1\,e^{-\alpha\,(i-1)} \ge \eta
  \;\Longrightarrow\;
  i-1 \le \frac{1}{\alpha}\log\!\frac{n_1}{\eta}
  \;\Longrightarrow\;
  \boxed{\;n_{\text{BNO}}(\eta)\;\sim\;\frac{1}{\alpha}\,\log\!\frac{n_1}{\eta}
  \;\approx\;\frac{1}{\alpha}\,\log\frac{1}{\eta}\;}
  \tag{$\star$}
\]
即\textbf{BNO 数随阈值按对数增长}: $\eta$ 减小一个数量级 $\to$ $n_{\text{BNO}}$
增加 $(\ln 10)/\alpha\approx 2.3/\alpha$ 个轨道。$\alpha$ 越大 (相关性越局域、
衰减越快), 同样 $\eta$ 下需要的 BNO 越少。

\subsection*{(3) 与幂律对照}
若 $\Delta\gamma$ 谱是幂律 $n_i\sim i^{-\beta}$ ($\beta>1$), 则
$n_i\ge\eta\Rightarrow i\le(n_1/\eta)^{1/\beta}$, 截断 BNO 数
$n_{\text{BNO}}\sim\eta^{-1/\beta}$ \emph{幂律增长}, 比 ($\star$) 慢很多
(指数 $-1/\beta$ vs $\log$)。实测 (见 \texttt{solution.py}) 在 H$_2$O/$6$-31G
上 $n_{\text{BNO}}$ vs $\eta$ 在 log-log 图接近直线、斜率约 $-0.2$, 既非纯对数
也非纯幂律 —— 因 $\alpha$ 不严格常数 (前面 BNO 衰减更快, 后面变慢)。
本题取 ($\star$) 作为\emph{主要} 渐近形式 (大 $i$ 时指数项主导), 它给出
\textbf{最稳健的工程结论}: $\eta$ 减半 $\to$ $n_{\text{BNO}}$ 仅增 $(\ln 2)/\alpha$
$\approx 0.7/\alpha$ 个轨道, 故 $\eta$ 越小 \emph{每多花一个 bath 轨道换来的精度
提升越小} —— 这正是 diminishing returns 的根源。

%=============================================================================
\section{Fragment size 与能量误差随 $\eta$ 的变化}
%=============================================================================
\subsection*{Fragment size}
每个 fragment cluster 的轨道数 (空间轨道)
\[
  n_{\text{cluster}}(\eta)
  = \underbrace{n_{\text{imp}}}_{\text{fragment}}
  + \underbrace{n_{\text{bath}}^{\text{DMET}}}_{\le n_{\text{imp}},\ \text{Schmidt}}
  + \underbrace{n_{\text{BNO,occ}}(\eta)+n_{\text{BNO,vir}}(\eta)}_{\text{MP2 bath, 截断于}\ \eta}.
\]
前两项与 $\eta$ 无关 (DMET minimal bath 固定), 故 cluster 随 $\eta$ 的增长完全由
BNO 项主导:
\[
  \boxed{\;n_{\text{cluster}}(\eta)\;\approx\;\text{const}+\frac{1}{\alpha}\big(\log(1/\eta_{\text{occ}})+\log(1/\eta_{\text{vir}})\big)\;}
  \quad (\eta_{\text{occ}}=\eta_{\text{vir}}=\eta\ \text{默认})。
\]
对应量子比特数 $n_{\text{qubit}}=2\,n_{\text{cluster}}$, 也按对数随 $\eta\!\downarrow$ 增长。
\textbf{对 NISQ 硬件的影响}: $\eta$ 减小 $\to$ $n_{\text{qubit}}$ 缓慢 (对数) 增长,
\emph{但} 量子电路深度 $\sim O(n_{\text{qubit}}^2)$, 故深度按 $(\log(1/\eta))^2$
增长 —— 这是 ``还能加多少 bath'' 的硬件上限。

\subsection*{能量误差}
设 ``无穷 bath'' ($\eta\to 0$) 极限下 EWF 能量为 $E_{\infty}^{\text{bath}}$,
\textbf{定义为最小 $\eta$ 下 EWF 的收敛能量} (实测: $\eta\le3\times10^{-6}$ 后
EWF 能量不再变)。注意 $E_{\infty}^{\text{bath}}$ 一般\emph{不严格等于} 全空间
CCSD —— 二者之差是 ``EWF 方法自身偏差'' (cluster 求解器 + 能量泛函近似), 与
bath 截断\emph{无关}。因此衡量 ``bath 截断误差'' $\epsilon_{\text{frag}}$ 应以
$E_{\infty}^{\text{bath}}$ 为参考, 而非 CCSD。

截掉占据 $<\eta$ 的 BNO 引入的误差大致\emph{正比于被截掉的占据之和}
(每条 BNO 贡献的相关能 $\propto n_i$):
\[
  \epsilon_{\text{frag}}(\eta)=|E_{\text{EWF}}(\eta)-E_{\infty}^{\text{bath}}|
  \;\sim\;\sum_{i>n_{\text{BNO}}(\eta)} n_i
  \;\sim\;\sum_{i>n_0} n_1\,e^{-\alpha(i-1)}
  \;\sim\; \frac{n_1}{\alpha}\,e^{-\alpha\, n_{\text{BNO}}}
  \;\sim\;\frac{n_1}{\alpha}\,\eta .
\]
即\textbf{bath 截断能量误差近似正比于 $\eta$}:
\[
  \boxed{\;\epsilon_{\text{frag}}(\eta)\;\propto\;\eta\;}
  \qquad(\text{指数衰减谱下}).
\]
\textbf{含义}: $\eta$ 减半 $\to$ 误差减半 (一阶), 但 $n_{\text{BNO}}$ 只增
$(\ln 2)/\alpha$ 个 —— \emph{精度按 $\eta$ 线性改善, 代价 (bath 规模) 按对数
增长}。这看起来 ``划算'', 但下一节会看到当 $\eta$ 进入 ``噪声/SQD-采样'' 量级
时, 边际收益被其他误差淹没。

%=============================================================================
\section{拐点与 diminishing returns}
%=============================================================================
\subsection*{边际分析}
定义两条 ``边际曲线'':
\begin{itemize}
  \item \textbf{边际成本}: 多保留一个 BNO 增加的计算代价
        $\dfrac{\partial\,\text{Cost}}{\partial n_{\text{BNO}}}
        \sim \dfrac{\partial\,n_{\text{cluster}}^3}{\partial n_{\text{BNO}}}
        \sim 3\,n_{\text{cluster}}^2$ (CCSD/SQD 对角化是 $O(n^3)$--$O(n^4)$)。
        因 $n_{\text{cluster}}$ 随 $\eta\!\downarrow$ \emph{增长}, 边际成本\emph{单调上升}。
  \item \textbf{边际收益}: 多保留一个 BNO 减少的误差
        $\dfrac{\partial\epsilon}{\partial n_{\text{BNO}}}\sim -n_{n_{\text{BNO}}+1}$
        $\propto -e^{-\alpha\,n_{\text{BNO}}}$, \emph{随} $\eta\!\downarrow$
        (即 $n_{\text{BNO}}\!\uparrow$) \emph{指数下降}。
\end{itemize}
\textbf{拐点} 出现在\emph{边际成本 $=$ 边际收益} 的 $\eta^\star$:
\[
  \underbrace{3\,n_{\text{cluster}}(\eta^\star)^2}_{\text{边际成本}}
  \;\sim\;
  \underbrace{e^{-\alpha\,n_{\text{BNO}}(\eta^\star)}}_{\text{边际收益}}
  \;\Longrightarrow\;
  \eta^\star \;\sim\; \text{const}\cdot n_{\text{cluster}}^2\,e^{-\alpha\,n_{\text{BNO}}}.
\]
更工程化的判据: 当 $\epsilon_{\text{frag}}(\eta)$ 降到与\emph{其他误差源}
(SQD 采样误差 $\epsilon_{\text{SQD}}$, 量子噪声 $\epsilon_{\text{noise}}$) \emph{同量级}
时, 再降 $\eta$ 已无意义:
\[
  \boxed{\;\eta^\star:\quad \epsilon_{\text{frag}}(\eta^\star)\;\approx\;\epsilon_{\text{SQD}}+\epsilon_{\text{noise}}\;}
  \tag{$\eta^\star$}
\]
这是 ``木桶效应'' 在 bath threshold 上的体现 (与第 9 问一致): 把
$\epsilon_{\text{frag}}$ 压到远小于其他误差, 总误差不变。

\subsection*{数值判据 (实测, H$_2$O/6-31G, EWF-CCSD solver)}
\sloppy
\texttt{solution.py} 在 H$_2$O/6-31G (13 空间轨道, 3 个原子 fragment) 上扫描
$\eta\in\{10^{-2},\ldots,10^{-6}\}$, 以最小 $\eta$ 下 EWF 能量
$E_{\infty}^{\text{bath}}=-76.11851$ Ha 为参考 (全空间 CCSD $=-76.11935$ Ha,
二者差 $\sim 8\times10^{-4}$ Ha 是 EWF 方法偏差, 与 bath 截断无关):

\begin{center}
\begin{tabular}{ccccccc}
\toprule
$\eta$ & $\bar n_{\text{BNO,occ}}$ & $\bar n_{\text{BNO,vir}}$
       & $\bar n_{\text{cluster}}$ & $\max n_{\text{cluster}}$
       & $E_{\text{EWF}}$ (Ha) & $\epsilon_{\text{frag}}$ (Ha) \\
\midrule
$10^{-2}$ & 0.0 & 0.0 & 3.7 & 7  & $-76.0229$ & $9.6\times10^{-2}$ \\
$10^{-3}$ & 0.0 & 1.0 & 4.7 & 10 & $-76.0903$ & $2.8\times10^{-2}$ \\
$10^{-4}$ & 0.7 & 3.3 & 7.7 & 13 & $-76.1095$ & $9.0\times10^{-3}$ \\
$10^{-5}$ & 2.0 & 4.7 & 10.3 & 13 & $-76.1163$ & $2.2\times10^{-3}$ \\
$10^{-6}$ & 2.7 & 6.0 & 11.7 & 13 & $-76.1185$ & $\to 0$ (收敛) \\
\bottomrule
\end{tabular}
\end{center}

\textbf{观察}:
\begin{itemize}
  \item BNO 占据数指数衰减拟合: 占据 (occ) $\alpha_{\text{occ}}\approx0.77$
        ($R^2{=}0.87$), 虚 (vir) $\alpha_{\text{vir}}\approx0.49$ ($R^2{=}0.97$),
        \emph{验证 (E) 的指数假设}; $\alpha_{\text{vir}}<\alpha_{\text{occ}}$ 说明
        虚 BNO 衰减更慢 (动态相关更离域到虚空间)。
  \item $\eta$ 从 $10^{-2}\to10^{-6}$ (4 个数量级), 平均 cluster 仅从 $3.7\to11.7$
        轨道 ($\sim 3\times$), \emph{对数增长} 与 ($\star$) 一致。
  \item $\epsilon_{\text{frag}}$ 从 $10^{-1}\to0$ Ha, 在
        $\eta\approx 10^{-4}$--$10^{-5}$ 处跨过 $10^{-2}\to10^{-3}$ 的 ``mHa 阈''
        —— \textbf{这正是拐点 $\eta^\star\!\sim\!10^{-4}$--$10^{-5}$}:
        再降 $\eta$ 到 $10^{-6}$, cluster 又加 1--2 条 BNO, 但 $\epsilon_{\text{frag}}$
        只从 $2{\times}10^{-3}\to0$, 进入 diminishing returns。
\end{itemize}

\subsection*{为什么 ``diminishing returns''?}
直观图景: MP2 BNO 的占据数谱是\emph{单调下降} 的, 前几条 BNO (大占据)
贡献绝大部分相关能 (Pareto 80/20); 后面的 BNO 占据 $\le\eta^\star$ 后,
其相关贡献 $\propto n_i$ 已\emph{指数小}, 而每多保留一条 BNO 都让 cluster
$n$ 加 1 $\to$ 对角化代价 $\sim n^3$ 加一截。当
\[
  \underbrace{\Delta\epsilon_{\text{frag}}\sim n_i}_{\text{边际收益}}
  \;<\;
  \underbrace{\text{常数}\cdot n_{\text{cluster}}^2}_{\text{边际代价 (CPU/GPU 时间)}}
\]
时, 进一步降 $\eta$ 是\emph{负 ROI}. 这就是为什么 EWF 默认
$\eta\sim 10^{-4}$--$10^{-5}$ (vayesta 默认 $10^{-5}$): 在典型弱相关分子上,
它已把 $\epsilon_{\text{frag}}$ 压到 mHa 以下, 与 CCSD 求解器自身精度相当,
再降 $\eta$ 只是 ``用更多算力换更小的尾部 BNO'', 不改善实际精度。

%=============================================================================
\section{特殊情况: STO-3G 基组的限制}
%=============================================================================
题目未限定基组, 但\textbf{STO-3G (7 个空间轨道) 对 H$_2$O 太小}, 会出现以下情况:
\begin{itemize}
  \item O 原子 fragment ($n_{\text{imp}}=5$): DMET minimal bath
        ($\le n_{\text{imp}}=5$) 已经几乎用完所有 7 个轨道, \emph{环境轨道为 0},
        故 O-fragment \textbf{没有 MP2 BNO 可加} (occup 数组为空);
  \item H 原子 fragment ($n_{\text{imp}}=1$): 仅 4 个 occ 环境 BNO + 1 个 vir,
        占据数从 $6{\times}10^{-4}$ 衰减到 $5{\times}10^{-7}$, 是少数能扫描的
        体系, 但 BNO 总数太少, $\eta$ 扫描曲线退化。
\end{itemize}
因此 \texttt{solution.py} 主扫描用 \textbf{6-31G (13 轨道)} 给出有意义的
$n_{\text{BNO}}$ vs $\eta$ 曲线 (环境足够大, 指数衰减可见);
同时提供 STO-3G 的诊断输出, 说明基组选择的物理含义。
这与 vayesta 官方 BNO 示例 (\texttt{examples/ewf/molecules/03-bno-threshold.py}
用 aug-cc-pVDZ) 一致 —— BNO 扫描需要足够大的基组才有意义。

%=============================================================================
\section*{代码与结果}
%=============================================================================
见 \texttt{solution.py}:
\begin{enumerate}
  \item H$_2$O/6-31G (主) + STO-3G (诊断) RHF, 报 $E_{\text{MF}}$;
  \item 全空间 CCSD 作参考 $E_{\infty}\approx E_{\text{CCSD}}$;
  \item vayesta EWF, bath\_options=\{bathtype="mp2", threshold=$\eta$\},
        扫描 $\eta\in\{10^{-6},...,10^{-2}\}$, 记录
        $n_{\text{BNO,occ/vir}}$, cluster 大小, $E_{\text{EWF}}$;
  \item 拟合 $n_i\sim n_1 e^{-\alpha(i-1)}$ 给出衰减常数 $\alpha$, 验证
        $n_{\text{BNO}}(\eta)\sim\tfrac{1}{\alpha}\log(1/\eta)$;
  \item 画 (a) BNO 占据数衰减 + 指数拟合;
        (b) fragment size vs $\eta$;
        (c) 能量误差 vs $\eta$ + 拐点标注
        $\to$ \texttt{bath\_threshold.png};
  \item 数值定位拐点 $\eta^\star$ (误差二阶差分极大点), 讨论 diminishing returns。
\end{enumerate}

\end{document}
